%%--------------------------------------------------------------------
%1--------------------------------------------------------------------
\vspace*{\fill}
\begin{glyph}{Hard (固)}{hard}\label{glyph:hard}
Hard.  Solid.
\end{glyph}

\begin{comps}
\comprtwocone & 囗 + 古\\\hline
\comprthreecone & Enclosure/Prison + Old (\ref{glyph:old})\\\hline
\end{comps}

\begin{remglyph}
\hooglig{Hard}, \remcom{old} \remcom{prison}.\\\hline
\end{remglyph}

\begin{prons}
\pronsrtwocone & \textbf{コ (KO)} & \textbf{かた (kata)}\\\hline
\pronsrthreecone & --- & --- \\\hline
\end{prons}

\begin{remprons}
\hooglig{A hard} \comon{copy} of the \comkun{catalogue}.\\\hline
\end{remprons}

%{\middel}
% &  \mbox{()}\\\hline
\begin{somme}
固い & \textit{hard; solid} & かた{\middel}い \mbox{(kata{\middel}i)}\\\hline
固体 & hard + body = \textit{a solid} & コ{\middel}タイ \mbox{(KO{\middel}TAI)}\\\hline
固定 & hard + fixed = \textit{fixed; settled} & コ{\middel}テイ \mbox{(KO{\middel}TEI)}\\\hline
\end{somme}

\begin{decomp}{7}
水&は&凍る&と&固体&に&なる。\\\hline
みず&は&こおる&と&コタイ&に&なる\\\hline
water&は&freeze&と&solid&に&to become\\\hline
\multicolumn{7}{|r|}{\textit{Water becomes solid when it freezes.}}\\\hline
\end{decomp}

\end{multicols}
\vhrulefill{2pt}
%%--------------------------------------------------------------------
%2--------------------------------------------------------------------
\begin{multicols}{2}
\begin{glyph}{Cool (冷)}{cool}\label{glyph:cool}
Cool.  Cold.\\\hline
Objects that are cool or cold to the touch (the idea of cold weather, however, is generally conveyed by the 漢字 in Entry \ref{glyph:cold}).
\end{glyph}

\begin{comps}
\comprtwocone & 冫 + 令 \\\hline
\comprthreecone & Ice + Command (\ref{glyph:command}) \\\hline
\end{comps}

\begin{remglyph}
\hooglig{Cool} \remcom{command} to make \remcom{ice}.\\\hline
\end{remglyph}

\begin{prons}
\pronsrtwocone & \textbf{レイ (REI)} & \textbf{つめ (tsume)}\\\hline
\pronsrthreecone & --- & ひ、さ (hi, sa) \\\hline
\end{prons}

\begin{remprons}
\hooglig{Cool} \comon{rain} with \ucomkun{heaps} of \comkun{stewed melon} in \ucomkun{summary}.\\\hline
\end{remprons}

%{\middel}
% &  \mbox{()}\\\hline
\begin{somme}
冷たい & \textit{cool} & つめ{\middel}たい \mbox{(tsume{\middel}tai)}\\\hline
冷気 & cool + spirit = \textit{chill; cold weather} & レイ{\middel}キ \mbox{(REI{\middel}KI)}\\\hline
\notyet{寒冷} & cold + cool = \textit{cold; chilly} & カン{\middel}レイ \mbox{(KAN{\middel}REI)}\\\hline
\end{somme}

\begin{decomp}{5}
今日&は&風&が&冷たい。\\\hline
きょう&は&かぜ&が&つめたい\\\hline
today&は&wind&が&chilly\\\hline
\multicolumn{5}{|r|}{\textit{The wind is cold today.}}\\\hline
\end{decomp}

\end{multicols}
\vspace*{\fill}
\pagebreak
%%--------------------------------------------------------------------
%3--------------------------------------------------------------------
\vspace*{\fill}
\begin{multicols}{2}
\begin{glyph}{Long Ago (昔)}{long_ago}\label{glyph:long_ago}
Long Ago.  Old Times.
\end{glyph}

\begin{comps}
\comprtwocone & 共 + 日 \\\hline
\comprthreecone & Together + Sun \\\hline
\end{comps}

\begin{remglyph}
\hooglig{Long ago}, the \remcom{sun} was together with \remcom{humanity}.\\\hline
\end{remglyph}

\begin{prons}
\pronsrtwocone & --- & \textbf{むかし (mukashi)}\\\hline
\pronsrthreecone & セキ、シャク (SEKI, SHAKU) & --- \\\hline
\end{prons}

\begin{remprons}
\hooglig{Long ago} they were \ucomon{sexy} \comkun{mukashi} koi in \ucomon{shucks}.\\\hline
{\warn}Mukashi Koi is a type of koi that has a metallic colour.\\\hline
\end{remprons}

%{\middel}
% &  \mbox{()}\\\hline
\begin{somme}
昔 & \textit{long ago} & むかし \mbox{(mukashi)}\\\hline
昔話 & long ago + speak = \textit{old story; folklore} & むかし{\middel}ばなし \mbox{(mukashi{\middel}banashi)}\\\hline
\end{somme}

\begin{decomp}{3}
昔&は&良かった。\\\hline
むかし&は&よかった\\\hline
olden days&は&good\\\hline
\multicolumn{3}{|r|}{\textit{I've seen better days.}}\\\hline
\end{decomp}

\end{multicols}
\vhrulefill{2pt}
%%--------------------------------------------------------------------
%4--------------------------------------------------------------------
\begin{multicols}{2}
\begin{glyph}{Inquiry (究)}{inquiry}\label{glyph:inquiry}
Inquiry, often with the sense of this being done thoroughly or extensively.
\end{glyph}

\begin{comps}
\comprtwocone & 穴 + 九\\\hline
\comprthreecone & Cave + Nine (Snail)\\\hline
\end{comps}

\begin{remglyph}
\hooglig{Inquire} about the \remcom{cave} \remcom{snail}.\\\hline
\end{remglyph}

\begin{prons}
\pronsrtwocone & \textbf{キュウ (KY\={U})} & \textbf{きわ (kiwa)}\\\hline
\pronsrthreecone & --- & --- \\\hline
\end{prons}

\begin{remprons}
\hooglig{Inquiring} about the \comkun{key-warden} will \comon{kill you}.\\\hline
\end{remprons}

%{\middel}
% &  \mbox{()}\\\hline
\begin{somme}
究める & \textit{investigate thoroughly; master} & きわ{\middel}める \mbox{(kiwa{\middel}meru)}\\\hline
\notyet{研究} & polish + inquiry = \textit{research} & ケン{\middel}キュウ \mbox{(KEN{\middel}KY\={U})}\\\hline
\end{somme}

\begin{decomp}{7}
彼&は&日本&文学&の&研究者&だ。\\\hline
かれ&は&ニホン&ブンガク&の&ケンキュウシャ&だ\\\hline
he&は&Japan&literature&の&researcher&be/is\\\hline
\multicolumn{7}{|r|}{\textit{He's a student of Japanese literature.}}\\\hline
\end{decomp}

\end{multicols}
\vspace*{\fill}
\pagebreak
%%--------------------------------------------------------------------
%5--------------------------------------------------------------------
\vspace*{\fill}
\begin{multicols}{2}
\begin{glyph}{Medicine (薬)}{medicine}\label{glyph:medicine}
Medicine.\\\hline
In a few compounds, this can mean chemical.
\end{glyph}

\begin{comps}
\comprtwocone & 楽 \\\hline
\comprthreecone & Music (\ref{glyph:music}) \\\hline
\end{comps}

\begin{remglyph}
\hooglig{Medicine} is a \remcom{melody} of \remcom{herbs}.\\\hline
\end{remglyph}

\begin{prons}
\pronsrtwocone & \textbf{ヤク (YAKU)} & \textbf{くすり (kusuri)}\\\hline
\rye{3}{The 訓読み is always voiced outside of the first position, as in the second below.}
\pronsrthreecone & --- & --- \\\hline
\end{prons}

\begin{remprons}
\hooglig{The medicine} for \comon{Jaco} consists of \comkun{couscous uric acid}.\\\hline
\end{remprons}

%{\middel}
% &  \mbox{()}\\\hline
\begin{somme}
薬 & \textit{medicine} & くすり \mbox{(kusuri)}\\\hline
目薬 & eye + medicine = \textit{eye drops} & め{\middel}ぐすり \mbox{(me{\middel}gusuri)}\\\hline
薬局 & medicine + office = \textit{pharmacy} & ヤッ{\middel}キョク \mbox{(YAK{\middel}KYOKU)}\\\hline
\end{somme}

\begin{decomp}{5}
良薬&は&口&に&苦し。\\\hline
リョウヤク&は&くち&に&にがし\\\hline
good medicine&は&mouth&に&bitter\\\hline
\multicolumn{5}{|r|}{\textit{A good medicine tastes bitter.}}\\\hline
\end{decomp}

\end{multicols}
\vhrulefill{2pt}
%%--------------------------------------------------------------------
%6--------------------------------------------------------------------
\begin{multicols}{2}
\begin{glyph}{Line (線)}{line}\label{glyph:line}
Line.
\end{glyph}

\begin{comps}
\comprtwocone & 糸 + 泉 \\\hline
\comprthreecone & Thread + Spring(s) (\ref{glyph:springs}) \\\hline
\end{comps}

\begin{remglyph}
\hooglig{Line} \remcom{threads} in the \remcom{spring}.\\\hline
\end{remglyph}

\begin{prons}
\pronsrtwocone & \textbf{セン (SEN)} & --- \\\hline
\pronsrthreecone & --- & --- \\\hline
\end{prons}

\begin{remprons}
\hooglig{Line} up with \comon{Central}.\\\hline
\end{remprons}

%{\middel}
% &  \mbox{()}\\\hline
\begin{somme}
本線 & main + line = \textit{main railway line} & ホン{\middel}セン \mbox{(HON{\middel}SEN)}\\\hline
前線 & before + line = \textit{front line; (weather) front} & ゼン{\middel}セン \mbox{(ZEN{\middel}SEN)}\\\hline
配線 & distribute + line = \textit{wiring} & ハイ{\middel}セン \mbox{(HAI{\middel}SEN)}\\\hline
\end{somme}

\begin{decomp}{5}
彼&は&前線&へ&やられた。\\\hline
かれ&は&ゼンセン&へ&やられた\\\hline
he&は&front line&へ&to dispatch\\\hline
\multicolumn{5}{|r|}{\textit{He was sent into combat.}}\\\hline
\end{decomp}

\end{multicols}
\vspace*{\fill}
\pagebreak
%%--------------------------------------------------------------------
%7--------------------------------------------------------------------
\vspace*{\fill}
\begin{multicols}{2}
\begin{glyph}{Face (顔)}{face}\label{glyph:face}
Face.
\end{glyph}

\begin{comps}
\comprtwocone & 彦 + 頁 \\\hline
\comprthreecone & Boy (archaic) / Lad + Big shell / Page / Face \\\hline
\end{comps}

\begin{remglyph}
\hooglig{The face} of the \remcom{lad} with \remcom{big shells} on this \remcom{page}.\\\hline
\end{remglyph}

\begin{prons}
\pronsrtwocone & \textbf{ガン (GAN)} & \textbf{かお (kao)}\\\hline
\rye{3}{The 訓読み is almost always voiced when not in the first position, as in the third compound.}
\pronsrthreecone & --- & --- \\\hline
\end{prons}

\begin{remprons}
\hooglig{Face} mask of \comon{gunpowder} and \comkun{kaolin}.\\\hline
\end{remprons}

%{\middel}
% &  \mbox{()}\\\hline
\begin{somme}
顔 & \textit{face} & かお \mbox{(kao)}\\\hline
洗顔 & wash + face = \textit{washing one's face} & セン{\middel}ガン \mbox{(SEN{\middel}GAN)}\\\hline
\notyet{横顔} & sideways + face = \textit{profile} & よこ{\middel}がお \mbox{(yoko{\middel}gao)}\\\hline
\end{somme}

\begin{decomp}{4}
顔&が&青い&よ。\\\hline
かお&が&あおい&よ\\\hline
face&が&blue (pale)&よ\\\hline
\multicolumn{4}{|r|}{\textit{You look pale.}}\\\hline
\end{decomp}

\end{multicols}
\vhrulefill{2pt}
%%--------------------------------------------------------------------
%8--------------------------------------------------------------------
\begin{multicols}{2}
\begin{glyph}{Tribe (族)}{tribe}\label{glyph:tribe}
Tribe.  Family.
\end{glyph}

\begin{comps}
\comprtwocone & 方 + ⺅ + 矢 \\\hline
\comprthreecone & Way/Direction + Person + Arrow\\\hline
\end{comps}

\begin{remglyph}
\hooglig{Tribes} \remcom{direct} their \remcom{people} like \remcom{arrows}.\\\hline
\end{remglyph}

\begin{prons}
\pronsrtwocone & \textbf{ゾク (ZOKU)} & --- \\\hline
\pronsrthreecone & --- & --- \\\hline
\end{prons}

\begin{remprons}
\hooglig{Tribe} \comon{Zock On!}\\\hline
{\warn}Zock On! is a single by Japanese hip hop group Teriyaki Boyz.\\\hline
\end{remprons}

%{\middel}
% &  \mbox{()}\\\hline
\begin{somme}
家族 & house + tribe = \textit{family} & カ{\middel}ゾク \mbox{(KA{\middel}ZOKU)}\\\hline
民族 & people + tribe = \textit{race; a people} & ミン{\middel}ゾク \mbox{(MIN{\middel}ZOKU)}\\\hline
\notyet{水族館} & water + tribe + manor = \textit{public aquarium} & スイ{\middel}ゾク{\middel}カン \mbox{(SUI{\middel}ZOKU{\middel}KAN)}\\\hline
\end{somme}

\begin{decomp}{5}
何&人&家族&です&か。\\\hline
なん&ニン&カゾク&です&か\\\hline
how many&people&family&be / is&か\\\hline
\multicolumn{5}{|r|}{\textit{How many people are there in your family.}}\\\hline
\end{decomp}

\end{multicols}
\vspace*{\fill}
\pagebreak
%%--------------------------------------------------------------------
%9--------------------------------------------------------------------
\vspace*{\fill}
\begin{multicols}{2}
\begin{glyph}{Teach (教)}{teach}\label{glyph:teach}
Teach.\\\hline
A religious element can be present in this 漢字, as shown in the second compound below.
\end{glyph}

\begin{comps}
\comprtwocone & 孝 + 攴攵 \\\hline
\comprthreecone & Filial Piety + Whip \\\hline
\end{comps}

\begin{remglyph}
\hooglig{Teach} with the \remcom{whip} to instil \remcom{filial piety}.\\\hline
\end{remglyph}

\begin{prons}
\pronsrtwocone & \textbf{キョウ (KY\={O})} & \textbf{おし (oshi)}\\\hline
\pronsrthreecone & --- & おそ (oso) \\\hline
\end{prons}

\begin{remprons}
\hooglig{Teach} \ucomkun{awesome} \comon{Kyoto} \comkun{osiers}.\\\hline
\end{remprons}

%{\middel}
% &  \mbox{()}\\\hline
\begin{somme}
教える & \textit{teach} & おし{\middel}える \mbox{(oshi{\middel}eru)}\\\hline
教会 & teach + meet = \textit{church} & キョウ{\middel}カイ \mbox{(KY\={O}{\middel}KAI)}\\\hline
\notyet{教科書} & teach + course + write = \textit{textbook; schoolbook} & キョウ{\middel}カ{\middel}ショ \mbox{(KY\={O}{\middel}KA{\middel}SHO)}\\\hline
\end{somme}

\begin{decomp}{5}
あの&教科書&は&もう&古い。\\\hline
あの&キョウカショ&は&もう&ふるい\\\hline
that&schoolbook&は&now&old\\\hline
\multicolumn{5}{|r|}{\textit{That textbook is out of date.}}\\\hline
\end{decomp}

\end{multicols}
\vhrulefill{2pt}
%%--------------------------------------------------------------------
%10-------------------------------------------------------------------
\begin{multicols}{2}
\begin{glyph}{Third (丙)}{third}\label{glyph:third}
Third.\\\hline
You will rarely meet this 漢字 on its own; its value for us is as a component in other more common characters.  It is, however, found in document of a more official nature, where it is used to indicate the ``third item'' or ``clause in a list''.
\end{glyph}

\begin{comps}
\comprtwocone & 一 + 内 \\\hline
\comprthreecone & One + Inside\\\hline
\end{comps}

\begin{remglyph}
\hooglig{Third} \remcom{one} \remcom{inside}.\\\hline
\end{remglyph}

\begin{prons}
\pronsrtwocone & \textbf{ヘイ (HEI)} & --- \\\hline
\pronsrthreecone & --- & --- \\\hline
\end{prons}

\begin{remprons}
\hooglig{Thirdly}, go \comon{haywire}.\\\hline
\end{remprons}

%{\middel}
% &  \mbox{()}\\\hline
\begin{somme}
丙 & \textit{third (in a series)} & ヘイ \mbox{(HEI)}\\\hline
\end{somme}

\begin{decomp}{11}
物事&の&第&３&番目&に&あたる&評価&が&丙&という。\\\hline
ものごと&の&ダイ&サン&ばんめ&に&あたる&ヒョウカ&が&ヘイ&という\\\hline
{\tiny thing}&{\tiny の}&{\tiny number}&{\tiny 3}&{\tiny position}&{\tiny に}&{\tiny relating to}&{\tiny rating}&{\tiny が}&{\tiny ヘイ}&{\tiny called}\\\hline
\multicolumn{11}{|r|}{\textit{The third grade or rank from the top is called `hei'.}}\\\hline
\end{decomp}

\end{multicols}
\vspace*{\fill}
\pagebreak
%%--------------------------------------------------------------------
%11-------------------------------------------------------------------
\vspace*{\fill}
\begin{multicols}{2}
\begin{glyph}{Loding (宿)}{lodging}\label{glyph:lodging}
Loding, from inns to hotels.
\end{glyph}

\begin{comps}
\comprtwocone & 宀 + 佰 \\\hline
\comprthreecone & Roof + 100 men \\\hline
\end{comps}

\begin{remglyph}
\hooglig{Lodging} with \remcom{100 men} under the \remcom{roof}.\\\hline
\end{remglyph}

\begin{prons}
\pronsrtwocone & \textbf{シュク (SHUKU)} & \textbf{やど (yado)}\\\hline
\pronsrthreecone & --- & --- \\\hline
\end{prons}

\begin{remprons}
\hooglig{The lodging} was \comon{shook} by \comkun{ya-dog}.\\\hline
\end{remprons}

%{\middel}
% &  \mbox{()}\\\hline
\begin{somme}
宿 & \textit{lodging; inn} & やど \mbox{(yado)}\\\hline
下宿 & lower + lodging = \textit{lodgings; boarding house} & ゲ{\middel}シュク \mbox{(GE{\middel}SHUKU)}\\\hline
\notyet{宿題} & lodging + topic = \textit{homework} & シュク{\middel}ダイ \mbox{(SHUKU{\middel}DAI)}\\\hline
\end{somme}

\begin{decomp}{3}
宿題&で&忙しい。\\\hline
シュクダイ&で&いそがしい\\\hline
homework&で&busy / occupied\\\hline
\multicolumn{3}{|r|}{\textit{I'm busy with my homework.}}\\\hline
\end{decomp}

\end{multicols}
\vhrulefill{2pt}
%%--------------------------------------------------------------------
%12-------------------------------------------------------------------
\begin{multicols}{2}
\begin{glyph}{Try (試)}{try}\label{glyph:try}
Try.
\end{glyph}

\begin{comps}
\comprtwocone & 言 + 式 \\\hline
\comprthreecone & Say + Type (\ref{glyph:type}) \\\hline
\end{comps}

\begin{remglyph}
\hooglig{Try} to this \remcom{type} of \remcom{saying}.\\\hline
\end{remglyph}

\begin{prons}
\pronsrtwocone & \textbf{シ (SHI)} & \textbf{こころ、ため (kokoro, tame)}\\\hline
\rye{3}{The 2 訓読み are verb stems for the first and second examples below.  Though they both mean to ``try'', 試す (ため{\middel}す/tame{\middel}su) has the added implication of \textit{testing something} so as to ascertain its value or validity.}
\pronsrthreecone & --- & --- \\\hline
\end{prons}

\begin{remprons}
\hooglig{Try} \comkun{taming} the \comkun{coco-roast} instead of \comon{shielding} against it.\\\hline
\end{remprons}

%{\middel}
% &  \mbox{()}\\\hline
\begin{somme}
\rye{3}{Note the mix of 音読み and 訓読み in the third compound below.}
試みる & \textit{try} & こころ{\middel}みる \mbox{(kokoro{\middel}miru)}\\\hline
試す & \textit{try; test} & ため{\middel}す \mbox{(tame{\middel}su)}\\\hline
試合 & try + match = \textit{game; match} & シあい \mbox{(SHI{\middel}ai)}\\\hline
\end{somme}

\begin{decomp}{7}
彼&は&種々&の&方法&を&試みた。\\\hline
かれ&は&シュジュ&の&ホウホウ&を&こころみた\\\hline
he&は&variety&の&way / method&を&to try\\\hline
\multicolumn{7}{|r|}{\textit{He tried many different methods.}}\\\hline
\end{decomp}

\end{multicols}
\vspace*{\fill}
\pagebreak
%%--------------------------------------------------------------------
%13-------------------------------------------------------------------
\vspace*{\fill}
\begin{multicols}{2}
\begin{glyph}{Anew (更)}{anew}\label{glyph:anew}
This 漢字 conveys a sense of doing things anew, be it through revision, reform or renewal, amongst others.\\\hline
A secondary meaning (expressed with the less common kunyomi)　has to do with staying up late, or the day growing late.
\end{glyph}

\begin{comps}
\comprtwocone & 日 + 一 + 丿乀 \\\hline
\comprthreecone & Sun/Day + One + Slash \\\hline
\end{comps}

\begin{remglyph}
\hooglig{Anew}, the \remcom{sun} \remcom{slashed} through its \remcom{one} and only path.\\\hline
\end{remglyph}

\begin{prons}
\pronsrtwocone & \textbf{コウ (K\={O})} & \textbf{さら (sara)}\\\hline
\pronsrthreecone & --- & ふ (fu) \\\hline
\end{prons}

\begin{remprons}
\hooglig{Anew}, the \ucomkun{foosball} \comon{caused} a \comkun{sarabande}.\\\hline
\end{remprons}

%{\middel}
% &  \mbox{()}\\\hline
\begin{somme}
更に & \textit{anew; furthermore} & さら{\middel}に \mbox{(sara{\middel}ni)}\\\hline
更生 & anew + life = \textit{rebirth; revival} & コウ{\middel}セイ \mbox{(K\={O}{\middel}SEI)}\\\hline
変更 & alter + anew = \textit{alteration; change} & ヘン{\middel}コウ \mbox{(HEN{\middel}K\={O})}\\\hline
\end{somme}

%\begin{decomp}{}
%\\\hline
%\\\hline
%\\\hline
%\multicolumn{}{|r|}{\textit{}}\\\hline
%\end{decomp}

\end{multicols}
\vhrulefill{2pt}
%%--------------------------------------------------------------------
%14-------------------------------------------------------------------
\begin{multicols}{2}
\begin{glyph}{System (系)}{system}\label{glyph:system}
System.\\\hline
As can be seen from the final compounds below, this 漢字 also appears in words related to lineage and genealogy.
\end{glyph}

\begin{comps}
\comprtwocone & 丿乀 + 糸 \\\hline
\comprthreecone & Slash (\#{4}) + Thread/Fabric\\\hline
\end{comps}

\begin{remglyph}
\hooglig{Our lineage} is \remcom{no} \remcom{fabrication}.\\\hline
\end{remglyph}

\begin{prons}
\pronsrtwocone & \textbf{ケイ (KEI)} & --- \\\hline
\pronsrthreecone & --- & --- \\\hline
\end{prons}

\begin{remprons}
\hooglig{A system} that \comon{caters} for all your needs.\\\hline
\end{remprons}

%{\middel}
% &  \mbox{()}\\\hline
\begin{somme}
体系 & body + system = \textit{system} & タイ{\middel}ケイ \mbox{(TAI{\middel}KEI)}\\\hline
系図 & system + diagram = \textit{family tree; genealogy} & ケイ{\middel}ズ \mbox{(KEI{\middel}ZU)}\\\hline 
日系 & sun (Japan) + system = \textit{of Japanese descent} & ニッ{\middel}ケイ \mbox{(NIK{\middel}KEI)}\\\hline
\end{somme}

%\begin{decomp}{}
%\\\hline
%\\\hline
%\\\hline
%\multicolumn{}{|r|}{\textit{}}\\\hline
%\end{decomp}

\end{multicols}
\vspace*{\fill}
\pagebreak
%%--------------------------------------------------------------------
%15-------------------------------------------------------------------
\vspace*{\fill}
\begin{multicols}{2}
\begin{glyph}{Depart (去)}{depart}\label{glyph:depart}
Depart.  Leave.
\end{glyph}

\begin{comps}
\comprtwocone & 厶 + 土 \\\hline
\comprthreecone & Self/Private + Earth/Soil \\\hline
\end{comps}

\begin{remglyph}
\hooglig{Depart} \remcom{privately} from the \remcom{Earth}.\\\hline
\end{remglyph}

\begin{prons}
\pronsrtwocone & \textbf{キョ (KYO)} & \textbf{さ (sa)}\\\hline
\pronsrthreecone & コ (KO) & --- \\\hline
\end{prons}

\begin{remprons}
\hooglig{The departing} \ucomon{cops} will \comon{keep your} \comkun{summaries}.\\\hline
\end{remprons}

%{\middel}
% &  \mbox{()}\\\hline
\begin{somme}
去る & \textit{depart; leave} & さ{\middel}る \mbox{(sa{\middel}ru)}\\\hline
去年 & depart + year = \textit{last year} & キョ{\middel}ネン \mbox{(KYO{\middel}NEN)}\\\hline
\notyet{飛び去る} & fly + depart = \textit{fly away} & と{\middel}び　さ{\middel}る \mbox{(to{\middel}bi sa{\middel}ru)}\\\hline
\end{somme}

%\begin{decomp}{}
%\\\hline
%\\\hline
%\\\hline
%\multicolumn{}{|r|}{\textit{}}\\\hline
%\end{decomp}

\end{multicols}
\vhrulefill{2pt}
%%--------------------------------------------------------------------
%16-------------------------------------------------------------------
\begin{multicols}{2}
\begin{glyph}{Nose (鼻)}{nose}\label{glyph:nose}
Nose.
\end{glyph}

\begin{comps}
\comprtwocone & 鼻 \\\hline
\comprthreecone & Nose\\\hline
\end{comps}

\begin{remglyph}
Part of the 漢字 radicals (\#{209}).\\\hline
\end{remglyph}

\begin{prons}
\pronsrtwocone & --- & \textbf{はな (hana)}\\\hline
\pronsrthreecone & ビ (BI) & --- \\\hline
\end{prons}

\begin{remprons}
\hooglig{Nose} \comkun{hung-ups} from snobs \ucomon{big-time}.\\\hline
\end{remprons}

%{\middel}
% &  \mbox{()}\\\hline
\begin{somme}
鼻 & \textit{nose} & はな \mbox{(hana)}\\\hline
鼻歌 & nose + song = \textit{humming} & はな{\middel}うた \mbox{(hana{\middel}uta)}\\\hline
鼻水 & nose + water = \textit{runny nose} & はな{\middel}みず \mbox{(hana{\middel}mizu)}\\\hline
\end{somme}

%\begin{decomp}{}
%\\\hline
%\\\hline
%\\\hline
%\multicolumn{}{|r|}{\textit{}}\\\hline
%\end{decomp}

\end{multicols}
\vspace*{\fill}
\pagebreak
%%--------------------------------------------------------------------
%17-------------------------------------------------------------------
\vspace*{\fill}
\begin{multicols}{2}
\begin{glyph}{Asia (亜)}{asia}\label{glyph:asia}
This 漢字 can stand for Asia in the same way that 日 does for Japan and 中 does for Chine.\\\hline
In a few other compounds, it can act like the English prefix `sub-'.\\\hline
It's an uncommon character, however, and will prove more useful as a component.
\end{glyph}

\begin{comps}
\comprtwocone & 口 + 一 + 丨 \\\hline
\comprthreecone & Mouth + One + Pole \\\hline
\end{comps}

\begin{remglyph}
\hooglig{The Asian} \remcom{one} \remcom{impaled} the \remcom{vampire}.\\\hline
\end{remglyph}

\begin{prons}
\pronsrtwocone & \textbf{ア (A)} & ---\\\hline
\pronsrthreecone & --- & --- \\\hline
\end{prons}

\begin{remprons}
\hooglig{In Asia}, things are always going \comon{upward}.\\\hline
\end{remprons}

%{\middel}
% &  \mbox{()}\\\hline
\begin{somme}
東亜 & East + Asia = \textit{Eastern Asia} & トウ{\middel}ア \mbox{(T\={O}{\middel}A)}\\\hline
\end{somme}

%\begin{decomp}{}
%\\\hline
%\\\hline
%\\\hline
%\multicolumn{}{|r|}{\textit{}}\\\hline
%\end{decomp}

\end{multicols}
\vhrulefill{2pt}
%%--------------------------------------------------------------------
%18-------------------------------------------------------------------
\begin{multicols}{2}
\begin{glyph}{Join (接)}{join}\label{glyph:join}
Join, Contact.
\end{glyph}

\begin{comps}
\comprtwocone & 手扌 + 立 + 女 \\\hline
\comprthreecone & Hands/Actions + Stand/Erect + Woman/Female \\\hline
\end{comps}

\begin{remglyph}
\hooglig{Contact} the one whose \remcom{actions} one cannot \remcom{withstand}.\\\hline
\end{remglyph}

\begin{prons}
\pronsrtwocone & \textbf{セツ (SETSU)} & --- \\\hline
\pronsrthreecone & --- & --- \\\hline
\end{prons}

\begin{remprons}
\hooglig{Join} \comon{sets}.\\\hline
\end{remprons}

%{\middel}
% &  \mbox{()}\\\hline
\begin{somme}
接続 & join + continue = \textit{joining; connection} & セツ{\middel}ゾク \mbox{(SETSU{\middel}ZOKU)}\\\hline
接着 & join + wear = \textit{adhesion} & セッ{\middel}チャク \mbox{(SET{\middel}CHAKU)}\\\hline
\notyet{面接} & mask + join = \textit{interview} & メン{\middel}セツ \mbox{(MEN{\middel}SETSU)}\\\hline
\end{somme}

%\begin{decomp}{}
%\\\hline
%\\\hline
%\\\hline
%\multicolumn{}{|r|}{\textit{}}\\\hline
%\end{decomp}

\end{multicols}
\vspace*{\fill}
\pagebreak
%%--------------------------------------------------------------------
%19-------------------------------------------------------------------
\vspace*{\fill}
\begin{multicols}{2}
\begin{glyph}{Opinion (説)}{opinion}\label{glyph:opinion}
Opinion, Theory.
\end{glyph}

\begin{comps}
\comprtwocone & 言 + 艹 + 兄 \\\hline
\comprthreecone & Speech + Grass/Herb/Plant + Elder Brother (\ref{glyph:elder_brother}) \\\hline
\end{comps}

\begin{remglyph}
\hooglig{Opinions} of your \remcom{elder brother} is like \remcom{speeches} of \remcom{grass stubble} to me.\\\hline
\end{remglyph}

\begin{prons}
\pronsrtwocone & \textbf{セツ (SETSU)} & --- \\\hline
\pronsrthreecone & ゼイ (ZEI) & と (to) \\\hline
\end{prons}

\begin{remprons}
\hooglig{Opinions} can \comon{set} you up fror \ucomon{zany}, \ucomkun{torrential} behaviour.\\\hline
\end{remprons}

%{\middel}
% &  \mbox{()}\\\hline
\begin{somme}
説明 & opinion + bright = \textit{explanation} & セツ{\middel}メイ \mbox{(SETSU{\middel}MEI)}\\\hline
伝説 & transmit + opinion = \textit{legend} & デン{\middel}セツ \mbox{(DEN{\middel}SETSU)}\\\hline
小説 & small + opinion = \textit{a novel} & ショウ{\middel}セツ \mbox{(SH\={O}{\middel}SETSU)}\\\hline
\end{somme}

%\begin{decomp}{}
%\\\hline
%\\\hline
%\\\hline
%\multicolumn{}{|r|}{\textit{}}\\\hline
%\end{decomp}

\end{multicols}
\vhrulefill{2pt}
%%--------------------------------------------------------------------
%20-------------------------------------------------------------------
\begin{multicols}{2}
\begin{glyph}{Make (作)}{make}\label{glyph:make}
Make.\\\hline
A secondary meaning relates to crops and harvesting.
\end{glyph}

\begin{comps}
\comprtwocone & 人亻 + 丿乀 + 丨 + 一 \\\hline
\comprthreecone & Man/Person + Slash + Pole + One \\\hline
\end{comps}

\begin{remglyph}
\hooglig{Make} a \remcom{pole} out of \remcom{one} \remcom{slashed} \remcom{person}.\\\hline
\end{remglyph}

\begin{prons}
\pronsrtwocone & \textbf{サク、サ (SAKU, SA)} & \textbf{つく (tsuku)}\\\hline
\rye{3}{サク is by far the more commonly used of these readings.  The kunyomi is usually voiced outside of the first position.}
\pronsrthreecone & --- & --- \\\hline
\end{prons}

\begin{remprons}
\hooglig{Make} me \comon{suck} on the \comon{subject} of \comkun{stew cooking}.\\\hline
\end{remprons}

%{\middel}
% &  \mbox{()}\\\hline
\begin{somme}
作る & \textit{make} & つく{\middel}る \mbox{(tsuku{\middel}ru)}\\\hline
作者 & make + individual = \textit{author; writer} & サク{\middel}シャ \mbox{(SAKU{\middel}SHA)}\\\hline
動作 & move + make = \textit{action; movement} & ドウ{\middel}サ \mbox{(D\={O}{\middel}SA)}\\\hline
\end{somme}

%\begin{decomp}{}
%\\\hline
%\\\hline
%\\\hline
%\multicolumn{}{|r|}{\textit{}}\\\hline
%\end{decomp}

\end{multicols}
\vspace*{\fill}
\pagebreak
%%--------------------------------------------------------------------
%21-------------------------------------------------------------------
\vspace*{\fill}
\begin{multicols}{2}
\begin{glyph}{Joint (節)}{joint}\label{glyph:joint}
Joint.\\\hline
A section of something.\\\hline
It is used in a variety of words, and can refer to everything from a `joint' between seasons (the second compound word provides an example), a knot in woord, of a melody in music, for instance.\\\hline
It can also convey a sense of moderation or economizing; think of this in terms of a narrowing joing that squeezes off the use of something such as the flow of money.
\end{glyph}

\begin{comps}
\comprtwocone & 竹ケ + 即 \\\hline
\comprthreecone & Bamboo + Immediate (\ref{glyph:immediate}) \\\hline
\end{comps}

\begin{remglyph}
\hooglig{Joint} the \remcom{bamboo} \remcom{immediately}.\\\hline
\end{remglyph}

\begin{prons}
\pronsrtwocone & \textbf{セツ (SETSU)} & \textbf{ふし (fushi)}\\\hline
\pronsrthreecone & セチ (SECHI) & --- \\\hline
\end{prons}

\begin{remprons}
\hooglig{Joint} \comkun{fusion} \comon{sets} \ucomon{sell cheap}.\\\hline
\end{remprons}

%{\middel}
% &  \mbox{()}\\\hline
\begin{somme}
節 & \textit{joint; melody; knot} & ふし \mbox{(fushi)}\\\hline
節分 & joint + part = \textit{the last day of winter} & セツ{\middel}ブン \mbox{(SETSU{\middel}BUN)}\\\hline
\notyet{調節} & tone + joint = \textit{adjustment; regulation} & チョウ{\middel}セツ \mbox{(CH\={O}{\middel}SETSU)}\\\hline
\end{somme}

%\begin{decomp}{}
%\\\hline
%\\\hline
%\\\hline
%\multicolumn{}{|r|}{\textit{}}\\\hline
%\end{decomp}

\end{multicols}
\vhrulefill{2pt}
%%--------------------------------------------------------------------
%22-------------------------------------------------------------------
\begin{multicols}{2}
\begin{glyph}{Rule (則)}{rule}\label{glyph:rule}
Rule, in the sense of law or regulation.
\end{glyph}

\begin{comps}
\comprtwocone & 貝 + 刂刁刀ク \\\hline
\comprthreecone & Shellfish + Knife/Sword \\\hline
\end{comps}

\begin{remglyph}
\hooglig{Rules} concerning \remcom{shellfish} \remcom{swords}.\\\hline
\end{remglyph}

\begin{prons}
\pronsrtwocone & \textbf{ソク (SOKU)} & --- \\\hline
\pronsrthreecone & --- & --- \\\hline
\end{prons}

\begin{remprons}
\hooglig{Rules} of \comon{sokkie}.\\\hline
\end{remprons}

%{\middel}
% &  \mbox{()}\\\hline
\begin{somme}
定則 & fixed + rule = \textit{established rule} & テイ{\middel}ソク \mbox{(TEI{\middel}SOKU)}\\\hline
原則 & original + rule = \textit{principle; general} & ゲン{\middel}ソク \mbox{(GEN{\middel}SOKU)}\\\hline
\notyet{法則} & law + rule = \textit{rule; law} & ホウ{\middel}ソク \mbox{(H\={O}{\middel}SOKU)}\\\hline
\end{somme}

%\begin{decomp}{}
%\\\hline
%\\\hline
%\\\hline
%\multicolumn{}{|r|}{\textit{}}\\\hline
%\end{decomp}

\end{multicols}
\vspace*{\fill}
\pagebreak
%%--------------------------------------------------------------------
%23-------------------------------------------------------------------
\vspace*{\fill}
\begin{multicols}{2}
\begin{glyph}{Now (今)}{now}\label{glyph:now}
Now.
\end{glyph}

\begin{comps}
\comprtwocone & 人 + 一 \\\hline
\comprthreecone & Man/Person + One \\\hline
\end{comps}

\begin{remglyph}
\hooglig{Now}, \remcom{man} has become \remcom{one}.\\\hline
\end{remglyph}

\begin{prons}
\pronsrtwocone & \textbf{コン (KON)} & \textbf{いま (ima)}\\\hline
\pronsrthreecone & キン (KIN) & --- \\\hline
\end{prons}

\begin{remprons}
\hooglig{Nowadays} \comkun{imams} \comon{kon} their \ucomon{kindred}.\\\hline
\end{remprons}

\begin{irrr}
今日 & now + sun (day) = \textit{today} & きょう \mbox{(ky\={o})}\\\hline
今朝 & now + morning = \textit{this morning} & けさ \mbox{(kesa)}\\\hline
今年 & now + year = \textit{this year} & こ{\middel}とし \mbox{(ko{\middel}toshi)}\\\hline
\rye{3}{We met the first two of these compounds in Entries 6 and 85, as the readings of both 漢字 were irregular.  The third we encounter here, as とし is the normal kunyomi of 年.  Also note how the addition of thei hiragana は in the second example below (which creates the most famous of all Japanese greetings) can change the readings of both 漢字 from the first example above.}
\end{irrr}

%{\middel}
% &  \mbox{()}\\\hline
\begin{somme}
今 & \textit{now} & いま \mbox{(ima)}\\\hline
今日は & now + sund (day) = \textit{Good morning / Afternoon; Hello!} & コン{\middel}ニチ{\middel}は \mbox{(KON{\middel}NICHI{\middel}wa)}\\\hline
今週 & now + week = \textit{this week} & コン{\middel}シュウ \mbox{(KON{\middel}SH\={U})}\\\hline
\end{somme}

%\begin{decomp}{}
%\\\hline
%\\\hline
%\\\hline
%\multicolumn{}{|r|}{\textit{}}\\\hline
%\end{decomp}

\end{multicols}
\vhrulefill{2pt}
%%--------------------------------------------------------------------
%24-------------------------------------------------------------------
\begin{multicols}{2}
\begin{glyph}{Eternal (永)}{eternal}\label{glyph:eternal}
Eternal.  Forever.\\\hline
Refer back to the 漢字 for ``water'' (Entry \#{61}) in order to see how the ``harpoon'' component of that character has become transformed in this one.
\end{glyph}

\begin{comps}
\comprtwocone & 丶 + 水 \\\hline
\comprthreecone & Dot + Water \\\hline
\end{comps}

\begin{remglyph}
\hooglig{Eternity} is not guaranteed by \remcom{dots} of \remcom{water} on your forehead.\\\hline
\end{remglyph}

\begin{prons}
\pronsrtwocone & \textbf{エイ (EI)} & \textbf{なが (naga)}\\\hline
\pronsrthreecone & --- & --- \\\hline
\end{prons}

\begin{remprons}
\hooglig{Eternally}, \comkun{Nagasaki} \comon{aims} for revenge.\\\hline
\end{remprons}

%{\middel}
% &  \mbox{()}\\\hline
\begin{somme}
永い & \textit{long (in terms of time)} & なが{\middel}い \mbox{(naga{\middel}i)}\\\hline
永遠 & eternal + far = \textit{eternity} & エイ{\middel}エン \mbox{(EI{\middel}EN)}\\\hline
\end{somme}

%\begin{decomp}{}
%\\\hline
%\\\hline
%\\\hline
%\multicolumn{}{|r|}{\textit{}}\\\hline
%\end{decomp}

\end{multicols}
\vspace*{\fill}
\pagebreak
%%--------------------------------------------------------------------
%25-------------------------------------------------------------------
\vspace*{\fill}
\begin{multicols}{2}
\begin{glyph}{Ice (氷)}{ice}\label{glyph:ice}
Ice.\\\hline
As you can see, this 漢字 is very similar visually to 永.  The only other thing to keep in mind here is that the stroke order of the `water' component is broken---the dot is written immediately after the initial harpoon strike.
\end{glyph}

\begin{comps}
\comprtwocone & 丶 + 水 \\\hline
\comprthreecone & Dot + Water \\\hline
\end{comps}

\begin{remglyph}
\hooglig{Icy} \remcom{water} \remcom{dots}.\\\hline
\end{remglyph}

\begin{prons}
\pronsrtwocone & \textbf{ヒョウ (HY\={O})} & \textbf{こおり (k\={o}ri)}\\\hline
\pronsrthreecone & --- & --- \\\hline
\end{prons}

\begin{remprons}
\hooglig{Icy} \comon{Hy\={o}rinmaru} eats \comkun{coriander}.\\\hline
\end{remprons}

%{\middel}
% &  \mbox{()}\\\hline
\begin{somme}
氷 & \textit{ice} & こおり \mbox{(k\={o}ri)}\\\hline
氷山 & ice + mountain = \textit{iceberg} & ヒョウ{\middel}ザン \mbox{(HY\={O}{\middel}ZAN)}\\\hline
氷結 & ice + bind = \textit{freeze over} & ヒョウ{\middel}ケツ \mbox{(HY\={O}{\middel}KETSU)}\\\hline
\end{somme}

%\begin{decomp}{}
%\\\hline
%\\\hline
%\\\hline
%\multicolumn{}{|r|}{\textit{}}\\\hline
%\end{decomp}

\end{multicols}
\vhrulefill{2pt}
\vspace*{\fill}
\pagebreak
%%--------------------------------------------------------------------
